\documentclass{article}
\usepackage{graphicx}
\usepackage{amsfonts}
\usepackage{amsmath}
\usepackage{amsthm}
\usepackage{tikz}

\theoremstyle{definition}
\newtheorem{theorem}{Theorem}[section]
\newtheorem{corollary}[theorem]{Corollary}
\newtheorem{proposition}[theorem]{Proposition}
\newtheorem{lemma}[theorem]{Lemma}
\newtheorem{definition}[theorem]{Definition}

\theoremstyle{remark}
\newtheorem{remark}[theorem]{Remark}
\newtheorem{criterion}[theorem]{Criterion}
\newtheorem{example}[theorem]{Example}
\newtheorem{claim}[theorem]{Claim}

\title{Complex Vector Space}
\author{Sarah Liu}
\date{January 2026}

\begin{document}

\maketitle
\section{Triangular Form}
\begin{definition}[Invariant]
    A subspace $W\subseteq V$ is invariant under $T$ if $T(W)\subseteq W$.
\end{definition}

\begin{proposition}
    Let $W=\text{span}\{\vec{x_1},\dots, \vec{x_k}\}\subseteq V$ is invariant if and only if $T(\vec{x_i})\subseteq W$ for all $i$.
\end{proposition}

\begin{definition}[Triangular Matrix]
    An $n\times n$ matrix is upper-triangular if 
    \[
    A=\begin{bmatrix}
        a_{11} & a_{12} & \dots & a_{1n}\\
        0 & a_{22} & \dots & a_{2n}\\
        \vdots & \vdots&&\vdots\\
        0 & 0 & \dots  & a_{nn}
    \end{bmatrix} \quad \text{if} \;a_{ij}=0\;\text{for}\; i>j
    \]
\end{definition}

\begin{definition}[Triangularizable]
    A linear transformation $T:V \rightarrow V$ is triangularizable if there exist a basis $\beta$ of $V$ such that $[T]^{\beta}_{\beta}$ is upper triangular.
\end{definition}

\begin{remark}
    If $T$ is diagonalizable, then it is triangularizable.
\end{remark}

\begin{theorem}
    Let $\beta=\{\vec{x_1},\dots,\vec{x_n}\}$ be a basis of $V$. Then $[T]_{\beta}^{\beta}$ is upper triangular if and only if $W_i=\text{span}\{\vec{x_1},\dots,\vec{x_i}\}$ are invariant under $T$ for any $i$.
\end{theorem}

\begin{proof}
    If $[T]^{\beta}_{\beta}$ is upper triangular, we could write it as
    \[
    [T]^{\beta}_{\beta}=\begin{bmatrix}
        a_{11} & a_{12} & \dots & a_{1n}\\
        0 & a_{22} & \dots & a_{2n}\\
        \vdots & \vdots&&\vdots\\
        0 & 0 & \dots  & a_{nn}
    \end{bmatrix}
    \]
    Then 
    \[
    \begin{aligned}
        T(\vec{x_1})&=a_{11}\vec{x_1}\in W_1\\
        T(\vec{x_2})&=a_{12}\vec{x_1}+a_{22}\vec{x_2}\in W_2\\
        T(\vec{x_3})&=a_{13}\vec{x_1}+a_{23}\vec{x_2}+a_{33}\vec{x_3}\in W_3\\
        \dots&\\
        T(\vec{x_n})&=a_{1n}\vec{x_1}+a_{2n}\vec{x_2}+\cdots +a_{nn}\vec{x_n}\in W_n
    \end{aligned}
    \]
    So $W_i$ is invariant under $T$.
    To prove the opposite, we can assume the basis and write it as the form above.
\end{proof}

\begin{remark}
    \[
    \vec{0}\in W_1\subseteq W_2\subseteq \cdots\subseteq W_n=V
    \]
\end{remark}

\begin{lemma}
    Let $T:V\rightarrow V$ be a linear transformation. If $W$ is invariant under $T$ and $\vec{x}\not\in W$, then $W+\vec{x}$ is invariant under $T$.
\end{lemma}

\begin{theorem}
    Let $V$ be a vector space under $\mathbb{C}$,.
    Let $T:V\rightarrow V$ be a linear transformation, then $T$ is triangularizable.
\end{theorem}

\begin{theorem}[Cayley-Hamiltonian Theorem]
    Let $T:V\rightarrow V$ be a linear transformation, and $V$ be a vector space over $\mathbb{C}$. 
    Let the characteristic polynomial be 
    \[
    p(t)=\chi_T(t)=a_n t^n+\cdots+a_1 t+a_0
    \]
    Then,
    \[
    p(T)=a_n T^n+\cdots+a_1 T+a_0=0
    \]
\end{theorem}

\section{Nilpotent Mapping}

\begin{definition}
    $N$ is nilpotent if $N^k=0$ for some $k\geq 1$.
\end{definition}

\begin{criterion}
    Let $N:V\rightarrow V$ be a linear transformation over finite dimensional vector space. Then $N$ is nilpotent if $N$ only has eigenvalue $\lambda=0$ with multiplicity $n$.
\end{criterion}

\begin{proof}
    Let $\vec{x}$ be an eigenvector, $N(\vec{x})=\lambda\vec{x}$, then $N^k(\vec{x})=\lambda^k\vec{x}=\vec{0}$, so $\lambda=0$.
\end{proof}

\begin{remark}
    If $N$ is nilpotent and $W$ is invariant under $N$, then $N|_W$ is nilpotent.
\end{remark}


\begin{definition}
    Let $\vec{x}\neq\vec{0}$.
    Let $N$ be nilpotent.
    Suppose $N^l(\vec{x})=\vec{0}$ and $N^{l-1}(\vec{x})\neq \vec{0}$.
    The cycle generated by $\vec{x}$ is 
    \[
    C\vec{x}=\{N^{l-1}(\vec{x}),N^{l-2}(\vec{x}),\dots, N(\vec{x}),\vec{x}\}
    \]
    The length of the cycle is $l$.
    The cyclic subspace generated by $\vec{x}$ is $C\vec{x}=\text{span}\{C\vec{x}\}$.
\end{definition}
\begin{remark}
    $C\vec{x}$ is an invariant subspace.
\end{remark}

\begin{proposition}
    \begin{itemize}
        \item $C\vec{x}$ is linearly independent
        \item $N^{l-1}(\vec{x})$ is an eigenvector
        \item $l\leq n=\dim(V)$

        If $l=n$, then $\alpha=C\vec{x}$ is a basis of $V$, and 
        \[
        [N]_{\alpha}^{\alpha}=\begin{bmatrix}
            0&1&0&\cdots &0\\
            0&0&1&\cdots &0\\
            \vdots&\vdots&\vdots&&\vdots\\
            0&0&0&\cdots &1\\
            0&0&0&\cdots &0\\
        \end{bmatrix}
        \]
    \end{itemize}
\end{proposition}

\begin{theorem}
    Let $C\vec{x_i}=\{N^{l_i-1}(\vec{x_i}),\cdots,N(\vec{x_i}),\vec{x_i}\}$ of length $l_i$ be the cycle generated by $\vec{x_i}$.
    If $\{N^{l_1-1}(\vec{x_1}),N^{l_2-1}(\vec{x_2}),\cdots, N^{l_k-1}(\vec{x_k})\}$ is linearly independent, then $C\vec{x_1}\cup\cdots\cup C\vec{x_k}$ is linearly independent.
\end{theorem}

\begin{definition}
    The cycles $C\vec{x_1}\cup\cdots\cup C\vec{x_k}$ are non-overlapping if $C\vec{x_1}\cup\cdots\cup C\vec{x_k}$ is linearly independent.
\end{definition}

\begin{definition}[Canonical Basis]
    Let $N$ be a nilpotent linear transformation. A basis $\beta$ is a canonical basis if $\beta=C\vec{x_1}\cup \cdots\cup C\vec{x_k}$ by some $\vec{x_1},\dots,\vec{x_t}$.
\end{definition}

\begin{theorem}
    Let $N:V\rightarrow V$ be a nilpotent transformation, then there exists a canonical basis.
\end{theorem}


\end{document}